\chapter{Grief}

\section*{Definition and Overview}
Grief is the natural emotional, physical, and psychological response to loss, most commonly the death of a loved one, but also to other losses such as the end of a relationship, loss of health, or loss of a job. Grief is not an illness; it is a normal and individual process that affects people differently. It can involve sadness, shock, anger, guilt, numbness, and a range of physical symptoms, and it typically comes in waves that ease over time. While most people adjust within a year or two, some experience prolonged or complicated grief that benefits from professional help. Understanding grief and allowing oneself to grieve in one's own way supports healing.

\section*{Epidemiology}
Grief is universal: virtually everyone experiences significant loss at some point in life. In Australia, around 160,000 people die each year, and each death affects an average of five close family members, so a substantial proportion of the population is grieving at any time. Most bereaved people experience significant distress that improves over months, but around 7--10\% develop prolonged grief disorder, in which intense grief persists beyond a year and impairs function. Grief also affects people who lose loved ones to suicide, accidents, or during the COVID-19 pandemic, where normal mourning rituals were disrupted.

\section*{Aetiology and Pathophysiology}
Grief is a response to loss rather than a disease, but it has biological, psychological, and social dimensions.

\begin{itemize}
    \item \textbf{Attachment:} humans form strong emotional bonds, and their loss triggers distress and yearning
    \item \textbf{Brain and body:} grief activates stress pathways, affecting sleep, appetite, immune function, and heart health
    \item \textbf{Psychological processes:} the mind works to make sense of the loss and adjust to a changed life
    \item \textbf{Cultural and social factors:} rituals, beliefs, and support networks shape how grief is expressed and experienced
    \item \textbf{Nature of the loss:} sudden, traumatic, or stigmatised deaths can complicate grief
    \item \textbf{Individual factors:} personality, past losses, and concurrent stresses influence the course of grief
\end{itemize}

\section*{Clinical Features}
Grief affects people in many ways, and no two people grieve alike.

\begin{itemize}
    \item Sadness, crying, and a sense of emptiness
    \item Shock, numbness, and disbelief, especially early on
    \item Yearning and longing for the person who died
    \item Anger, guilt, and regret
    \item Anxiety, fear, and a sense of insecurity
    \item Difficulty concentrating and poor memory
    \item Sleep disturbance, loss of appetite, and fatigue
    \item Physical symptoms, including chest tightness, headache, and a lowered resistance to illness
    \item Withdrawal from social activities
    \item Waves of intense emotion that come and go, often triggered by reminders
\end{itemize}

\section*{Differential Diagnosis}
Grief can resemble, or coexist with, mental health conditions that need treatment.

\begin{itemize}
    \item Clinical depression: persistent low mood with loss of interest, feelings of worthlessness, and poor function lasting most of the time for more than two weeks
    \item Prolonged grief disorder: intense grief that persists beyond a year and severely impairs daily life
    \item Anxiety and panic disorders
    \item Post-traumatic stress disorder, when the death was traumatic
    \item Complicated grief after a suicide or violent death
    \item Physical illness causing fatigue and weight loss
    \item Adjustment disorder in response to a non-death loss
\end{itemize}

\section*{When to Seek Medical Care}
See a GP if grief is overwhelming, persists beyond a year without improvement, or is accompanied by depression, thoughts of self-harm, or an inability to function at work or home. Seeking help for grief is a sign of strength, not weakness.

\textbf{Call 000 (or local emergency services) for:}
\begin{itemize}
    \item Thoughts of harming yourself or ending your life, or plans to do so
    \item Inability to care for yourself or severe withdrawal from the world
    \item Any crisis where you feel you cannot cope
\end{itemize}
For immediate support, Lifeline (13 11 14), GriefLine (1300 845 745), and Beyond Blue (1300 22 4636) are available around the clock.

\section*{Diagnosis and Investigations}
Grief is assessed through sensitive conversation rather than tests.

\begin{itemize}
    \item \textbf{Clinical interview:} exploring the nature of the loss, the relationship, and the person's responses
    \item \textbf{Assessment of duration and impact:} whether grief is interfering with work, relationships, and self-care
    \item \textbf{Screening for depression and anxiety:} standard questionnaires help identify coexisting conditions
    \item \textbf{Risk assessment:} asking about self-harm and suicidal thoughts is routine and important
    \item \textbf{Review of supports:} understanding the person's social network and coping resources
    \item \textbf{Physical assessment:} checking for physical symptoms that may need attention
\end{itemize}

\section*{Management}
Most grief resolves with time, support, and self-care; professional help is available when needed.

\begin{itemize}
    \item \textbf{Allowing grief:} giving yourself permission to feel and express emotions in your own way and time
    \item \textbf{Support:} talking with family, friends, and others who understand, including bereavement support groups
    \item \textbf{Rituals and remembrance:} funerals, memorials, and personal rituals help process the loss
    \item \textbf{Self-care:} maintaining sleep, eating, exercise, and routine
    \item \textbf{Counselling:} bereavement counselling and grief therapy help people process the loss
    \item \textbf{Cognitive behavioural therapy:} effective for complicated grief
    \item \textbf{Medication:} used only if depression or severe anxiety coexists, not for grief itself
    \item \textbf{Time:} recognising that grief eases gradually, and that there is no set timetable
\end{itemize}

\section*{Complications}
\begin{itemize}
    \item Prolonged grief disorder, with persistent yearning, emptiness, and disability
    \item Clinical depression and anxiety disorders
    \item Sleep disorders and fatigue
    \item Increased risk of physical illness, including heart disease, in the first months after loss
    \item Substance use as a way of coping
    \item Relationship strain and social isolation
    \item Work impairment and financial stress
    \item In severe cases, self-harm and suicide
\end{itemize}

\section*{Prognosis and Outcomes}
Most people adjust to loss within one to two years, with grief becoming less intense and less constant over time. The "waves" of grief gradually become less frequent and less overwhelming, and most people find a new balance while maintaining a continuing bond with the person they lost. About 7--10\% develop prolonged grief disorder, which responds well to specific therapy such as prolonged grief therapy or cognitive behavioural therapy. With support and, when needed, professional help, most people emerge from grief with renewed capacity for life and relationships.

\section*{Prevention and Self-Care}
\begin{itemize}
    \item Give yourself permission to grieve in your own way and at your own pace
    \item Talk about the person you lost and share memories with others
    \item Accept help from family and friends, and ask for practical support
    \item Maintain basic routines: regular meals, sleep, and gentle exercise
    \item Limit alcohol and avoid using drugs to numb the pain
    \item Join a bereavement support group to connect with others who understand
    \item Plan for anniversaries and significant dates, which can be hard
    \item Be patient with yourself, and avoid comparing your grief to others'
    \item Seek professional help early if grief is not easing or is affecting your ability to function
\end{itemize}

\section*{Sources and Further Reading}
The following reputable sources provide further information for healthcare students and patients seeking reliable, current advice about grief.

\begin{itemize}
    \item GriefLine. \textit{Support and counselling for grief and loss.} https://www.griefline.org.au/
    \item Australian Centre for Grief and Bereavement. \textit{Resources and services for bereavement.} https://www.grief.org.au/
    \item Lifeline Australia. \textit{Grief and loss support.} https://www.lifeline.org.au/
    \item Healthdirect Australia. \textit{Grief and loss.} https://www.healthdirect.gov.au/
    \item NHS. \textit{Grief after bereavement.} https://www.nhs.uk/
    \item American Psychological Association. \textit{Grief: coping with the loss of a loved one.} https://www.apa.org/
\end{itemize}
